\documentclass{article} % For LaTeX2e
\usepackage{iclr2022_conference,times}
% Optional math commands from https://github.com/goodfeli/dlbook_notation.
\input{math_commands.tex}
% Submission type
\newcommand{\apsname}{Yield Curve Jacobian}
% Group number
\newcommand{\gpnumber}{7}
\usepackage{hyperref}
\usepackage{url}
\usepackage{graphicx}
\usepackage{booktabs}
\usepackage{amsmath}
\usepackage{amssymb}
\usepackage{xcolor}
% Project title
\title{Leading-Order Degeneracy of a Decoder-Tangent Penalty:\\
A Constraint Ablation for Latent Yield-Curve Dynamics}
\author{Daniel Deng \\
Student\# 1011132141 \\
\texttt{ddaniel.deng@mail.utoronto.ca} \\
\And
Dylan Setyono \\
Student\# 1011116668 \\
\texttt{dylan.setyono@mail.utoronto.ca} \\
}
\newcommand{\fix}{\marginpar{FIX}}
\newcommand{\new}{\marginpar{NEW}}
\iclrfinalcopy
\begin{document}
\maketitle
\begin{abstract}
We study multi-horizon forecasting of the U.S.\ Treasury yield curve with a
two-stage geometric pipeline: a Student-$t$ conditional VAE learns a
low-dimensional curve representation, and a latent drift model evolves that state
under a persistence-residual decode, so that zero drift recovers a random walk.
On a shared frozen Stage~A model we compare three soft-training configurations---%
unconstrained, a decoder-Jacobian tangent penalty, and a bundle of curve-shape
regularizers---against persistence.
Our main contribution is analytic. We show that an exact tangent projector
annihilates the leading-order decoded latent move. In the joint limit of an exact
projector and a first-order decoder linearization, the penalty is constant in the
dynamics parameters. The implemented penalty is not in that limit: ridge
regularization creates both a reconstruction-residual cross-term and an anisotropic
shrinkage term, while decoder curvature supplies an additional training signal. We
derive these surviving terms and show that the familiar $\varepsilon/4$ spectral
bound applies only to the quadratic shrinkage component, not to the complete
penalty. Thus the construction loses its direct first-order tangent signal but is
not generally inert.
The empirical results are reported as point
estimates from a single training run. Adding the tangent penalty changes test RMSE
by at most $7\times10^{-5}$ at the trained horizons and reverses direction at the
extrapolated horizon. Both drift variants have errors close to persistence at
$h{=}1,5$ and improve on it at $h{=}21$ by roughly $1.3\%$ in point estimate;
persistence is better at
$h{=}63$. The configuration carrying the regularizer bundle diverges under
multi-step rollout. We perform no paired significance testing and vary no seeds, so
we make no claim that any of these differences is reproducible. These figures are
also conditional on a training panel containing the disclosed 2002--2006
interpolation defect. The contribution is a
characterization of the leading-order degeneracy and surviving signals of a
decoder-Jacobian tangent penalty, together with a constraint ablation on a shared
representation with a random walk as the null.
% Main body page count
----Total Pages: \pageref{last_page}
\end{abstract}
\section{Introduction}
The yield curve sits at the center of fixed-income finance: it prices bonds,
discounts future cash flows, and encodes the market's view of growth and monetary
policy. Modeling how the whole curve moves over time is what lets us forecast,
hedge, and stress-test rate portfolios. The task is more tractable than it looks,
because the curve is highly redundant. Since \citet{LittermanScheinkman1991},
three factors (level, slope, and curvature) have been known to explain almost all
of its variation. Classical models build on this in different ways:
Nelson--Siegel \citep{NelsonSiegel1987} and its dynamic Diebold--Li form
\citep{DieboldLi2006} impose fixed factor shapes, affine and HJM models
\citep{DuffieKan1996, HeathJarrowMorton1992} add explicit no-arbitrage structure,
and PCA captures the leading linear factors. These models are robust and
interpretable, but their linear, parametric form limits how well they track the
nonlinear way real curves deform across regimes.
Deep generative models offer more flexibility, but risk producing curves that
collapse to a few modes or violate no-arbitrage out of sample.
\citet{NoArbVAEYieldCurves} address this with a two-stage, physics-informed design:
a Student-$t$ conditional VAE with dynamic level injection learns a heavy-tailed
curve representation that separates macroeconomic shape from the absolute rate
level, and the latent state is then evolved by a continuous-time neural SDE
penalized by a no-arbitrage PDE. We take this pipeline as our starting point and
reproduce a deliberately simplified version of it: a single currency (USD), a
discrete tenor grid, and a first-order residual from which the Hessian convexity
term is omitted.
Alongside their main model, \citet{NoArbVAEYieldCurves} discuss a decoder-Jacobian
construction that ties the curve's admissible moves to the local geometry of the
learned representation, using the Jacobian as a volatility basis together with a
nonlinear re-encoding $D(E(\cdot))$. We do not implement that benchmark. What we
study instead is a construction suggested by the same geometric intuition and
substantially cheaper: a soft penalty that measures how far a forecast lies from
the affine tangent space at the decode of the current latent state, using a
ridge-regularized tangent map. We call it a linearized tangent-penalty surrogate to
keep the distinction from their proposal explicit, and our findings are about our
construction only.
Our main experiment is an ablation on a frozen Stage~A model. We train the latent
dynamics three ways---unconstrained, with the tangent penalty, and with a bundle of
curve-shape regularizers---so that comparisons are made on the same representation.
We evaluate on U.S.\ Treasury constant-maturity yields from FRED at horizons of one
day to one quarter, scoring point-forecast accuracy alongside re-encoding distance,
the off-tangent component of each forecast move, and forward-curve smoothness.
Critically, we benchmark against a random-walk (persistence) forecast. Persistence
is the appropriate null for daily yields and is not reported by
\citet{NoArbVAEYieldCurves}, whose headline figures concern reconstruction rather
than multi-step forecasting; including it changes which claims the evidence can
support.
Our central result is analytic and does not depend on any run. The decoded latent
move lies in $\operatorname{span}(J)$ to first order, so an exact tangent projector
annihilates its leading-order contribution. For the implemented ridge map, however,
the loss also contains a term coupling the latent move to the fixed Stage~A
reconstruction residual, as well as a quadratic shrinkage term; for the nonlinear
decoder, curvature creates a further signal. The $\varepsilon/4$ spectral bound
controls only the quadratic shrinkage term. Consequently, the construction has a
leading-order degeneracy, but the implemented penalty is not identically constant
and can still shape the dynamics through ridge leakage and decoder curvature.
The empirical findings are consistent with that analysis and are reported as point
estimates from a single training run per configuration. The persistence-residual
drift model is within $5\times10^{-5}$--$7\times10^{-4}$ of a random walk at daily
and weekly horizons, improves on it at $h{=}21$ by roughly $1.3\%$, and loses at the
extrapolated $h{=}63$. Adding the tangent penalty changes test RMSE by at most
$7\times10^{-5}$ at the trained horizons and reverses direction at $h{=}63$; the
geometry diagnostics likewise do not separate the two configurations in a consistent
direction. The regularizer bundle is the unstable configuration: worse than the
unconstrained model at every horizon, diverging from $0.041$ at $h{=}1$ to $7.68$ at
$h{=}63$.
Three limitations bound all of this and we state them at the outset. First, we varied
no seeds and ran no paired significance tests, so we make no claim that
sub-$10^{-4}$ differences are reproducible; observing a small difference is
consistent with weak surviving signals but does not establish their cause. Second,
Stage~A test reconstruction RMSE is $0.659$, so real curves sit far from
$D(E(y))$ and absolute re-encoding distance closely tracks the Stage~A error at the
persistence anchor. This makes projection-based accuracy gains uncertain, but it
does not constitute a mathematical bound on what a projection method could achieve.
Third, the training panel contains a four-year interpolated block caused by the
30-year-series discontinuation; all empirical findings are conditional on the
resulting Stage~A representation.
\section{Method}
\label{sec:method}
\subsection{Problem setup and notation}
Let $x_t \in \mathbb{R}^{N}$ denote the vector of constant-maturity Treasury
yields observed on business day $t$, where the $N=11$ maturities (tenors) are
\begin{equation*}
\tau \in \{1\mathrm{M}, 3\mathrm{M}, 6\mathrm{M}, 1\mathrm{Y}, 2\mathrm{Y},
3\mathrm{Y}, 5\mathrm{Y}, 7\mathrm{Y}, 10\mathrm{Y}, 20\mathrm{Y},
30\mathrm{Y}\},
\end{equation*}
with $\tau$ measured in years. We treat the observed panel
$\{x_t\}_{t=1}^{T}$ as discrete samples of a continuously evolving curve and seek
to model its dynamics so as to respect the well-documented low-dimensional
structure of the yield curve and to accommodate heavy-tailed yield innovations.
The pipeline is organized in two stages. \textbf{Stage~A} learns a low-dimensional
latent representation $z_t \in \mathbb{R}^{d}$ of the curve. \textbf{Stage~B} learns
the dynamics of $z_t$ and decodes forecasts back to curves, with optional soft
training penalties. We set $d=5$: the classical finding that approximately $99\%$ of
yield-curve variation is spanned by three factors motivates a low-dimensional
representation, and two additional dimensions give the decoder headroom for shape
variation the linear factors do not capture.
\subsection{Data and preprocessing}
\label{sec:data}
We obtain daily constant-maturity yields for the eleven tenors above from the
Federal Reserve Economic Data (FRED) repository. All results come from a single
panel and a single training run. The download is configured to start 2001-07-01, but
the binding constraint is the 1-month series (DGS1MO), whose first observation is
2001-07-31, so after incomplete rows are dropped the panel begins there. The frozen
scorecard does not record the day on which the panel was retrieved; the results were
committed on 2026-07-09 and a prior retrieval ran through 2026-06-25, which brackets
the end of the panel between those two dates and gives $T$ between $6{,}498$ and
$6{,}508$ business days. Nothing in what follows depends on which day within that
window the panel was pulled: two additional weeks at the end of the sample move the
$70\%$ cut by about a week.
The cleaning procedure has a consequence we state explicitly because it bears on how
the results should be read. We drop any day on which any tenor is missing, then
reindex the remaining panel to a regular business-day calendar and impute the
resulting gaps by time interpolation followed by forward/back-filling. The 30-year
series (DGS30) is discontinued between 2002-02-18 and 2006-02-09, so those days are
first removed by the completeness filter and then reinstated by the business-day
reindex and filled by interpolation: $1{,}037$ business days, roughly $16\%$ of the
panel, are linear interpolants rather than observations. The same mechanism creates
interpolated rows on market holidays, which the business-day calendar includes and
FRED does not report. Because interpolation is applied before the chronological
split and uses \texttt{limit\_direction="both"}, imputed values adjacent to a split
boundary can draw on observations from the following split.
The mitigating fact is that the interpolated 2002--2006 window falls entirely inside
the training split. Under the $70/15/15$ chronological partition, training runs from
2001-07-31 to early January 2019 (2019-01-03 or 2019-01-14, depending on the
retrieval day), validation from there to late September or mid-October 2022, and the
test split from autumn 2022 to the end of the panel in mid-2026. The interpolated
block therefore contaminates the learned representation and the scaling statistics,
but not the held-out evaluation on which every reported number is computed. We regard
a rebuild of the panel---beginning in 2006, or dropping the 30-year tenor---as the
correct fix and did not perform it; the results below should be read as conditional
on a Stage~A model trained on a panel with this defect.

Stage~B forecasts require a 21-business-day lookback, so evaluation begins 21
business days after the start of the test split. The held-out evaluation sample is
approximately $954$ forecast windows at $h{=}1$, falling to approximately $892$ at
$h{=}63$. Every figure and table below is scored on that sample, autumn 2022 to
mid-2026, and none of it overlaps the training period.
The cleaned panel is partitioned \emph{chronologically} into training ($70\%$),
validation ($15\%$), and test ($15\%$) sets; no shuffling is performed. We
standardize each tenor using robust statistics estimated \emph{on the training split
only},
\begin{equation}
  \tilde{x} = \frac{x - \operatorname{median}(x)}{\operatorname{IQR}(x)},
\end{equation}
and the fitted statistics are applied unchanged to the validation and test splits.
All reported metrics are in these scaled units.
We then apply a level/shape decomposition (``LevelScript''), \emph{after} robust
scaling: writing $\tilde{x}_t$ for the scaled curve, we take the level to be its
1-year component, $\ell_t = \tilde{x}_t^{(\mathrm{1Y})}$, and the shape to be
$s_t = \tilde{x}_t - \ell_t \mathbf{1}$. Because each tenor carries its own IQR,
subtracting the scaled 1-year component does not correspond to removing a parallel
shift of the underlying yield curve, and this ordering differs from the procedure of
\citet{NoArbVAEYieldCurves}, who decompose before scaling. We report the
decomposition as implemented and make no claim that it separates level from shape in
an economic sense; it supplies the decoder with a conditioning variable that carries
most of the curve's persistence.
Throughout, $D_\theta^{(s)}(z,\ell)$ denotes the decoded scaled shape and
$D_\theta(z,\ell)=D_\theta^{(s)}(z,\ell)+\ell\mathbf{1}$ the corresponding full
scaled curve; we write $y_t=\tilde{x}_t$ for that scaled full curve in the
forecasting equations. When a level argument is suppressed below, it is fixed at
the last level available at the forecast origin. Likewise, $E_\phi(x)$ is
shorthand for applying the encoder to the level/shape pair derived from the scaled
curve $x$.
\subsection{Stage A: latent representation learning}
\label{subsec:stageA}
Stage~A is a conditional variational autoencoder over scaled shape curves,
conditioned on the level. The approximate posterior is diagonal Gaussian,
\begin{equation}
  q_\phi(z \mid s, \ell) = \mathcal{N}\!\big(\mu_\phi(s,\ell),\,
  \operatorname{diag}\sigma_\phi^2(s,\ell)\big),
  \qquad p(z) = \mathcal{N}(0, I_d),
\end{equation}
and training maximizes the $\beta$-weighted evidence lower bound
\begin{equation}
  \mathcal{L}_{\mathrm{ELBO}}
  = \mathbb{E}_{q_\phi(z\mid s,\ell)}\big[\log p_\theta(s \mid z, \ell)\big]
  - \beta\, \mathrm{KL}\!\big(q_\phi(z \mid s,\ell)\,\|\,p(z)\big),
\end{equation}
with the expectation estimated by a single reparameterized sample
$z = \mu_\phi + \sigma_\phi \odot \epsilon$, $\epsilon \sim \mathcal{N}(0,I_d)$,
and $\beta = 0.03$.
The \textbf{Student-$t$ CVAE}, our Stage~A model, augments this with a per-tenor
Student-$t$ observation model,
\begin{equation}
  p_\theta(s_i \mid z, \ell) = t_{\nu}\!\left(s_i;\, \hat{s}_i,\, a_i\right),
  \qquad \hat{s}_i = D_\theta^{(s)}(z, \ell)_i,
\end{equation}
with per-tenor scale $a_i$ from a dedicated head and a learned degrees-of-freedom
parameter $\nu = \mathrm{softplus}(\nu_{\mathrm{raw}}) + 2.01$, constrained above
$2$ to keep the variance finite. The reconstruction term is the exact closed-form
Student-$t$ negative log-likelihood, whose polynomial tails make reconstruction
robust to large yield moves; the Gaussian model is the $\nu\to\infty$ limit. On
completion we export the deterministic latent means $z_t=\mu_\phi(s_t,\ell_t)$ for
all splits as the input series for Stage~B.
Stage~A uses latent dimension $d{=}5$, hidden width $256$, and $200$ epochs. On the
test set, reconstruction RMSE is $0.659$ (MAE $0.523$). This does not propagate
directly into the forecasts: the persistence-residual parameterization of
Section~\ref{subsec:stageB} decodes a \emph{difference} of two curves, so the static
reconstruction bias cancels to first order, which is why forecast RMSE at $h{=}1$
($\approx 0.029$) is an order of magnitude below the reconstruction error. It does,
however, govern the magnitude of the geometry diagnostics
(Section~\ref{sec:analysis-reencoding}) and the size of the residual term in
Section~\ref{sec:analysis-inert}. We do not report a rank-matched linear
reconstruction baseline, so we cannot separate the share of this error attributable
to the rank constraint from the share attributable to the decoder.
\subsection{Stage B: latent dynamics}
\label{subsec:stageB}
Using the frozen Stage~A encoder, the curve history maps to a latent trajectory
$\{z_t\}$. Write $\mathbf{z}_t=(z_{t-L+1},\dots,z_t)$ for the last $L=10$ latent
states, on which both networks are conditioned. Let $\psi$ denote the trainable
Stage~B parameters; the frozen decoder retains parameters $\theta$. The state is
modeled as
$dz_t=\mu_{\mathbb{P},\psi}(\mathbf{z}_t)\,dt+
\sigma_\psi(\mathbf{z}_t)\,dW_t^{\mathbb{P}}$,
with $\mu_{\mathbb{Q},\psi}=\mu_{\mathbb{P},\psi}-\sigma_\psi\lambda_\psi$.
Below we suppress $\psi$ and abbreviate $\mu_{\mathbb{P}}(\mathbf{z}_t)$ as
$\mu_{\mathbb{P}}(z_t)$ wherever neither dependence is at issue.
\paragraph{Persistence-residual forecasting.} Daily curves are close to a random
walk, so the last observed curve is a very strong baseline. We therefore
parameterize the $h$-step forecast as a correction on top of persistence,
\begin{equation}
  \hat{y}_{t+h} = y_t + D_\theta(z_{t+h},\ell_t) - D_\theta(z_t,\ell_t),
  \label{eq:persres}
\end{equation}
where $z_{t+h}$ is obtained by rolling the drift forward with step $\Delta t = 1$.
With zero drift this reduces exactly to persistence, so the model only has to learn
the departure from the random walk.
\paragraph{Training objective.} Each training batch draws a horizon
$h\sim\mathrm{Unif}\{1,5,21\}$. Stage~B minimizes a one-step latent term plus a
curve-space term on the decoded (persistence-residual) yields at that horizon,
\begin{equation}
  \mathcal{L}_{\mathrm{fit}}
  = w_{\mathrm{lat}}\,\big\|\hat{z}_{t+1}-z_{t+1}\big\|^2
  + w_{\mathrm{curve}}\,\big\|\hat{y}_{t+h}-y_{t+h}\big\|^2,
  \qquad \hat{z}_{t+1}=z_t+\mu_{\mathbb{P}}(z_t)\Delta t,
\end{equation}
with $w_{\mathrm{lat}}=0.05$ and $w_{\mathrm{curve}}=10$, both mean-squared errors.
The latent term is always one-step while the curve term follows the sampled horizon.
At the reported $h{=}1$ curve RMSE of $0.029$ the curve term is
$10\times(0.029)^2\approx 8.5\times10^{-3}$, so $\mathcal{L}_{\mathrm{fit}}$ is of
order $10^{-2}$ near convergence; we use this implied scale, rather than a logged
value, as the reference against which penalty magnitudes are compared in
Section~\ref{sec:analysis-inert}.
The best checkpoint is selected on validation curve-RMSE at $h{=}1$; no test data
enter model selection. To avoid lookahead, the LevelScript decoder uses the last
known level rather than the realized future level. Stage~B uses $400$ epochs,
learning rate $3\times10^{-4}$, and a cosine learning-rate schedule.
\paragraph{Role of the diffusion term.} Both Eq.~\eqref{eq:persres} and
$\mathcal{L}_{\mathrm{fit}}$ propagate the drift only, so the reported point
forecasts are deterministic drift rollouts---a plug-in approximation, not the exact
conditional mean of a nonlinear SDE---rather than sampled paths. The diffusion
$\sigma$ enters the objective only through the residual of
Section~\ref{subsec:constraints}: in the \texttt{sde\_only} and
\texttt{sde\_jacobian} configurations no term of the objective depends on $\sigma$,
and the diffusion is therefore not trained. Those two configurations are drift
networks, and we refer to them as such.
\subsection{Soft training penalties}
\label{subsec:constraints}
Both penalties are evaluated on the same persistence-residual forecast that the
curve-fitting loss scores, $\hat y = y_t + D_\theta(\hat z) - D_\theta(z_t)$, at the
horizon sampled for that batch, with horizon gating as described below. This is
worth stating precisely, because Section~\ref{sec:analysis-inert} turns on it: the
parameter-dependent part of $\hat y$ is the decoded latent move
$D_\theta(\hat z) - D_\theta(z_t)$, and the additive term $y_t$ is data.
\paragraph{Curve-shape regularizer bundle.} This configuration adds three terms. The
first is $\mathbb{E}[\mathcal{R}^2]$ with
$\mathcal{R}=-\partial_\tau P+\nabla_z P^{\top}\mu_{\mathbb{Q}}-rP$, in the form of
the first-order term of a pricing PDE, accompanied by discount-monotonicity and
forward-smoothness penalties. We deliberately do not call this a no-arbitrage
residual, for four reasons that together make the label indefensible. (i)~The
quantities entering $P=\exp(-y\tau)$ are the decoder's outputs, which are in
robust-scaled units rather than decimals, and $r$ is the scaled shortest-tenor
yield. Centered, dimensionless quantities cannot be interpreted as annualized
continuously compounded rates. Although a discount factor above one is possible
under genuinely negative rates, here its value is driven by preprocessing rather
than by a pricing convention. (ii)~FRED constant-maturity yields are par yields,
and zero-coupon prices cannot be recovered from them by $P=e^{-y\tau}$ without
bootstrapping, which we do not perform. (iii)~$\partial_\tau P$ is implemented as
$-yP$, which holds only if $y$ is constant in maturity; the
$\tau\,\partial_\tau y$ term is dropped, as is the Hessian convexity term of the
full PDE. (iv)~The learned drift uses one business day as its time unit, whereas
$\tau$ and $r$ enter the residual in annual units, with no factor of $252$ reconciling
the two. The bundle is therefore a heuristic regularizer on a transform of the
scaled curve and on the latent dynamics. Our findings about it carry no implication
for correctly specified no-arbitrage penalization. The bundle is applied at $h{=}1$
only.
\paragraph{Decoder-Jacobian tangent penalty.} The penalty
$w\big\|\hat y-\Pi_{\mathcal{M}}(\hat y)\big\|^2$ is applied at
$h\in\{1,5,21\}$ with weight $w=0.3$ and a $40$-epoch linear warmup. The tangent map
is ridge-regularized,
\begin{equation}
  \Pi_{\mathcal{M}}(\hat y)=D_\theta(z_t)+J(J^{\top}J+\varepsilon I)^{-1}J^{\top}
  \big(\hat y-D_\theta(z_t)\big),
  \qquad J = J(z_t)=\frac{\partial D_\theta}{\partial z}(z_t),
  \label{eq:ridgemap}
\end{equation}
with $\varepsilon = 10^{-5}$. Three properties matter for
Section~\ref{sec:analysis-inert}: for $\varepsilon>0$ the map is not idempotent, so
it is only an approximate projector; it targets the affine tangent space at
$D_\theta(z_t)$, a local linearization of the decoder image rather than that image
itself; and $J$ is evaluated at the encoder output $z_t$ with a frozen decoder, so
it is a constant with respect to the trainable Stage~B parameters $\psi$.
\subsection{Ablation design, hyperparameters, and reproducibility}
\label{sec:setup}
The ablation on a shared frozen Stage~A model compares Stage~B configurations that
share the same architecture and differ only in soft training penalties:
\begin{itemize}
  \item \texttt{sde\_only}: no penalties;
  \item \texttt{sde\_jacobian}: decoder-Jacobian tangent penalty at
        $h\in\{1,5,21\}$;
  \item \texttt{sde\_pde}: the curve-shape regularizer bundle, applied at $h{=}1$.
\end{itemize}
A fourth configuration combining both penalties was trained. It returned metrics
numerically identical to \texttt{sde\_pde} across all eight scored quantities at all
four horizons, which the training code does not predict, since the tangent penalty
carries a nonzero warmup-scaled weight from the first epoch. Distinguishing a
checkpoint-selection artifact from a configuration fault requires the per-epoch
training histories, which were not retained. We therefore exclude that configuration
rather than interpret it.
All reported Stage~B numbers are \emph{soft} forecasts: penalties act during training
only. A hard-projection variant applying the map of Eq.~\eqref{eq:ridgemap} at
inference exists in the codebase but is not used in any result reported here.
Locked hyperparameters are $d=5$, hidden width $256$, KL weight $0.03$, $200$
Stage~A epochs, and $400$ Stage~B epochs with a cosine schedule.
\paragraph{Reproducibility.} All reported runs use a single seed ($42$), applied via
\texttt{torch.manual\allowbreak\_seed} and
\texttt{numpy.random\allowbreak.seed}. Deterministic kernels
were not enabled and training ran on a cloud GPU, so runs are not bit-reproducible.
We did not vary the seed and we did not repeat any configuration under a controlled
protocol, so we have no variance estimate and we make no claim about whether the
sub-$10^{-4}$ differences reported below would survive a rerun. Per-variant
best-checkpoint epochs are written to the Stage~B history files, which were not
retained; no claim in this paper depends on them.
\subsection{Baseline}
We benchmark against \textbf{persistence} ($\hat{y}_{t+h}=y_t$), the dominant
short-horizon reference for daily curves. We also implemented rank-3 PCA--VAR and
Nelson--Siegel--Svensson/Diebold--Li baselines, but exclude their results because a
fault in the rolling refit duplicated overlapping observations in the estimation
window. Without corrected runs, those errors are not interpretable as estimates of
either model class's performance.
\subsection{Evaluation protocol}
We evaluate at horizons $h\in\{1,5,21,63\}$ business days; $h{=}63$ lies outside
the training horizon set and is treated as an extrapolation stress test. We report
mean-curve RMSE for accuracy. For geometry we report three quantities: the
re-encoding distance $\|\hat{y}-D_\theta(E_\phi(\hat{y}))\|$; the \emph{re-encoding
gain}, the same quantity minus its value at the persistence anchor $y_t$, which
removes the shared persistence-anchor residual and is negative when the forecast has
a smaller re-encoding residual than persistence; and a tangent-move residual
measuring the part of the decoded forecast move outside
$\operatorname{span}(J(z_t))$. We also report forward-curve smoothness. Note that
$D(E(y))$ is not the nearest point of the decoder image to $y$---the encoder is not
a metric projection---so these are re-encoding diagnostics, not distances to a
manifold.
\section{Results}
\label{sec:results}
All metrics are computed on the held-out \textbf{test} split in scaled yield units,
from a single training run per configuration. Lower is better throughout. No
confidence intervals or significance tests accompany these numbers; differences
should be read as point estimates.
\subsection{Forecast accuracy}
\label{sec:results-rmse}
Table~\ref{tab:curve-rmse} reports test curve RMSE across horizons. At short
horizons ($h{=}1,5$), persistence attains the lowest error, consistent with
near-random-walk behavior of daily Treasury yields. The two drift variants track it
closely.
At $h{=}21$---within the Stage~B training horizon set---both drift variants have
lower error than persistence ($0.1221$ and $0.1222$ vs.\ $0.1238$, a margin of about
$1.3\%$). The two differ from each other by $5\times10^{-5}$, and by at most
$7\times10^{-5}$ at any trained horizon; at $h{=}63$ the ordering reverses. We do not
interpret differences of this size in either direction.
The configuration carrying the regularizer bundle is already worse at $h{=}1$ and
deteriorates rapidly thereafter, from $0.041$ at $h{=}1$ to $7.68$ at $h{=}63$.
\begin{table}[t]
  \centering
  \caption{Test-set curve RMSE (scaled yields), single run per configuration. Bold
  marks the lowest entry in each column. Differences between \texttt{sde\_only} and
  \texttt{sde\_jacobian} are at most $7\times10^{-5}$ at the trained horizons and are
  not accompanied by any inferential test.}
  \label{tab:curve-rmse}
  \begin{tabular}{@{}lcccc@{}}
    \toprule
    Model & $h{=}1$ & $h{=}5$ & $h{=}21$ & $h{=}63$ \\
    \midrule
    Persistence
      & $\mathbf{0.02912}$ & $\mathbf{0.06233}$ & $0.12375$ & $\mathbf{0.20239}$ \\
    \midrule
    \texttt{sde\_only}
      & $0.02918$ & $0.06301$ & $0.12217$ & $0.23654$ \\
    \texttt{sde\_jacobian}
      & $0.02917$ & $0.06294$ & $\mathbf{0.12212}$ & $0.24637$ \\
    \texttt{sde\_pde}
      & $0.04101$ & $0.19675$ & $1.75664$ & $7.67525$ \\
    \bottomrule
  \end{tabular}
\end{table}
\subsection{Geometry diagnostics}
\label{sec:results-geometry}
Tables~\ref{tab:geom-h1}--\ref{tab:geom-h21} report the geometry and shape
diagnostics at the trained horizons.
\paragraph{Tangent-move residual.} This metric measures the off-tangent component of
the \emph{nonlinear} decoded latent move at $z_t$. It is not the quantity the
training penalty evaluates, which acts on the persistence-residual forecast and is
leading-order degenerate for the reason given in
Section~\ref{sec:analysis-inert}.
\texttt{sde\_jacobian} is nominally lower at all three horizons, by margins of
$1\times10^{-6}$, $1.3\times10^{-5}$ and $5.4\times10^{-5}$; with one run per
configuration we cannot attribute these to the penalty.
\paragraph{Re-encoding distance and gain.} Absolute re-encoding distance remains
large ($\approx0.65$) for all drift forecasts and closely tracks the Stage~A
re-encoding error at the persistence anchor, so it does not discriminate between
configurations. The
re-encoding gain, which removes that constant, does not order the two drift variants
consistently: \texttt{sde\_jacobian} is nominally better at $h{=}1$ and $h{=}5$
($0.00077$ vs.\ $0.00090$; $0.00341$ vs.\ $0.00397$) and nominally worse at $h{=}21$
($-0.00652$ vs.\ $-0.00744$, where more negative is better). We report both
directions rather than the favorable subset.
\paragraph{Regularizer contrast.} \texttt{sde\_pde} reduces absolute re-encoding
distance at short horizons but at the cost of large RMSE, and at $h{=}21$ degrades
sharply on every diagnostic, with a tangent residual of $0.85$. The stability of the
other two configurations is associated with the \emph{absence} of the bundle rather
than the presence of the tangent penalty.
\begin{table}[t]
  \centering
  \caption{Ablation at $h{=}1$ (test). Lower is better for RMSE, re-encoding
  distance, tangent residual and smoothness; for re-encoding gain, more negative is
  better.}
  \label{tab:geom-h1}
  \setlength{\tabcolsep}{4pt}
  \begin{tabular}{@{}lccccc@{}}
    \toprule
    Model & RMSE & Re-enc.\ dist. & Re-enc.\ gain & Tangent res. & Fwd.\ smooth. \\
    \midrule
    \texttt{sde\_only}
      & $0.02918$ & $0.6526$ & $0.00090$ & $2.9\!\times\!10^{-5}$ & $0.0375$ \\
    \texttt{sde\_jacobian}
      & $0.02917$ & $0.6525$ & $0.00077$ & $2.8\!\times\!10^{-5}$ & $0.0375$ \\
    \texttt{sde\_pde}
      & $0.04101$ & $0.6333$ & $-0.01838$ & $0.00330$ & $0.0487$ \\
    \bottomrule
  \end{tabular}
\end{table}
\begin{table}[t]
  \centering
  \caption{Ablation at $h{=}5$ (test).}
  \label{tab:geom-h5}
  \setlength{\tabcolsep}{4pt}
  \begin{tabular}{@{}lccccc@{}}
    \toprule
    Model & RMSE & Re-enc.\ dist. & Re-enc.\ gain & Tangent res. & Fwd.\ smooth. \\
    \midrule
    \texttt{sde\_only}
      & $0.06301$ & $0.6572$ & $0.00397$ & $0.000338$ & $0.0379$ \\
    \texttt{sde\_jacobian}
      & $0.06294$ & $0.6566$ & $0.00341$ & $0.000325$ & $0.0379$ \\
    \texttt{sde\_pde}
      & $0.19675$ & $0.5101$ & $-0.14385$ & $0.04935$ & $0.1060$ \\
    \bottomrule
  \end{tabular}
\end{table}
\begin{table}[t]
  \centering
  \caption{Ablation at $h{=}21$ (test). The regularizer bundle collapses; the two
  drift variants remain stable and are not ordered consistently across diagnostics.}
  \label{tab:geom-h21}
  \setlength{\tabcolsep}{4pt}
  \begin{tabular}{@{}lccccc@{}}
    \toprule
    Model & RMSE & Re-enc.\ dist. & Re-enc.\ gain & Tangent res. & Fwd.\ smooth. \\
    \midrule
    \texttt{sde\_only}
      & $0.12217$ & $0.6479$ & $-0.00744$ & $0.00180$ & $0.0379$ \\
    \texttt{sde\_jacobian}
      & $0.12212$ & $0.6489$ & $-0.00652$ & $0.00175$ & $0.0375$ \\
    \texttt{sde\_pde}
      & $1.75664$ & $1.2570$ & $0.61702$ & $0.84791$ & $3.0489$ \\
    \bottomrule
  \end{tabular}
\end{table}
\begin{figure}[t]
  \centering
  \includegraphics[width=0.48\linewidth]{stage_b_sde_only_rmse_vs_horizon.png}
  \hfill
  \includegraphics[width=0.48\linewidth]{stage_b_sde_pde_rmse_vs_horizon.png}
  \caption{Test curve RMSE against forecast horizon for \texttt{sde\_only} (left)
  and \texttt{sde\_pde} (right). The unconstrained drift degrades smoothly with
  horizon; the configuration carrying the regularizer bundle diverges beyond the
  horizon at which the bundle is applied. Note the difference in vertical scale.}
  \label{fig:horizon}
\end{figure}
\subsection{Summary of findings}
\label{sec:results-summary}
\begin{enumerate}
  \item \textbf{Accuracy.} Forecasts from the tangent-penalized and unconstrained
  drift variants differ by at most $7\times10^{-5}$ at the trained horizons, with the
  ordering reversing at $h{=}63$. Both have lower error than persistence at $h{=}21$
  by roughly $1.3\%$ in point estimate.
  \item \textbf{Geometry.} The two are not ordered consistently across the geometry
  diagnostics, and no diagnostic separates them by a margin we can attribute to the
  penalty with one run.
  \item \textbf{Regularizer bundle.} The configuration carrying it is unstable under
  multi-step evaluation. The design does not identify which of its three terms is
  responsible.
  \item \textbf{Baseline.} Persistence is the lowest-error forecast at every
  horizon except $h{=}21$.
\end{enumerate}
\section{Analysis}
\label{sec:analysis}
\subsection{Persistence as the dominant short-horizon reference}
\label{sec:analysis-persistence}
At $h{=}1$ and $h{=}5$, persistence attains the lowest curve RMSE ($0.02912$ and
$0.06233$). This is expected for daily Treasury yields: day-to-day changes are small
relative to the level of the curve, so ``tomorrow equals today'' is a strong
baseline. The drift variants sit within $5\times10^{-5}$--$7\times10^{-4}$ of
persistence at these horizons. On this panel, the informative comparison is
therefore whether adding the geometric term changes the forecast, not whether either
drift variant clears persistence by a practically meaningful margin.
The medium horizon $h{=}21$ is more informative, and both drift variants have lower
error than persistence there ($0.12212$ / $0.12217$ vs.\ $0.12375$). The margin is
about $1.3\%$, it is a point estimate from a single run, and it is not accompanied by
a paired test (Section~\ref{sec:analysis-limitations}). We treat $h{=}21$ as the
horizon where learned multi-step drift could plausibly matter, and $h{=}1,5$ as
near-random-walk regimes where matching persistence is the bar.
\subsection{Observed forecast changes are small}
\label{sec:analysis-jac-vs-only}
On the reported run, \texttt{sde\_jacobian} has marginally lower error than
\texttt{sde\_only} at every trained horizon:
\begin{center}
\begin{tabular}{@{}lccc@{}}
  \toprule
  Horizon & \texttt{sde\_jacobian} & \texttt{sde\_only} & Difference \\
  \midrule
  $h{=}1$  & $0.029167$ & $0.029175$ & $8\times10^{-6}$ \\
  $h{=}5$  & $0.062944$ & $0.063010$ & $6.6\times10^{-5}$ \\
  $h{=}21$ & $0.122123$ & $0.122173$ & $5.0\times10^{-5}$ \\
  $h{=}63$ & $0.246372$ & $0.236540$ & $-9.8\times10^{-3}$ \\
  \bottomrule
\end{tabular}
\end{center}
The differences at trained horizons are in the fourth and fifth decimal places,
amounting to at most $0.1\%$ of the error, and the ordering reverses at $h{=}63$. We
ran one seed per configuration and no paired test, so we cannot say whether these
differences are reproducible, and we do not claim either that the penalty improves
accuracy or that it is provably neutral in a statistical sense. What we can say is
that the observed change at trained horizons is small relative to the total forecast
error. The next section characterizes why the leading-order signal vanishes while
also identifying the terms through which the implemented penalty can still train.
\subsection{Leading-order degeneracy and surviving training signals}
\label{sec:analysis-inert}
Write $r_t = y_t - D_\theta(z_t)$ for the Stage~A reconstruction residual at the
forecast origin and
$m_\psi = D_\theta(\hat z_\psi) - D_\theta(z_t)$ for the decoded move controlled by
the Stage~B parameters $\psi$. With $A = J^{\top}J$,
$P_\varepsilon = J(A + \varepsilon I)^{-1}J^{\top}$, and
$Q_\varepsilon=I-P_\varepsilon$, the implemented penalty is
\begin{equation}
  \mathcal{L}_{\mathrm{tan}}(\psi)
  =w\big\lVert Q_\varepsilon\big(r_t+m_\psi\big)\big\rVert^2 .
  \label{eq:penalty-actual}
\end{equation}
Stage~A is frozen and $y_t$ is data, so $r_t$ does not depend on the dynamics
parameters. Because $J$ is evaluated at the encoder output $z_t$ with a frozen
decoder, $J$, $A$, $P_\varepsilon$ and $Q_\varepsilon$ are likewise constants with
respect to $\psi$; every expansion below treats them as fixed and differentiates only
through $m_\psi$. The projected component $Q_\varepsilon r_t$ is therefore constant
in $\psi$, but its interaction with the decoded move need not be.
\paragraph{Exact projection and a linearized move.} First replace the nonlinear
move by $J\delta_\psi$, where $\delta_\psi=\hat z_\psi-z_t$, and define
$B_\varepsilon=Q_\varepsilon J$. Expanding the loss gives
\begin{equation}
\begin{split}
  \mathcal{L}_{\mathrm{tan}}^{\mathrm{lin}}(\psi)
  ={}&w\lVert Q_\varepsilon r_t\rVert^2
  +2w\,\delta_\psi^{\top}B_\varepsilon^{\top}Q_\varepsilon r_t\\
  &+w\,\delta_\psi^{\top}B_\varepsilon^{\top}B_\varepsilon\delta_\psi .
  \label{eq:linearized-expansion}
\end{split}
\end{equation}
Let $P_0=JJ^+$, with $J^+$ the Moore--Penrose pseudoinverse, be the exact
orthogonal projector onto
$\operatorname{span}(J)$. Then $Q_0J=0$, so $B_0=0$ and both
dynamics-dependent terms in Eq.~\eqref{eq:linearized-expansion} vanish. Thus an
exact projector applied to a first-order decoded move supplies no training signal to
$\psi$, even when $r_t\ne0$. This is the precise degeneracy result.
\paragraph{Finite-ridge leakage.} For $\varepsilon>0$,
\begin{equation}
  B_\varepsilon=Q_\varepsilon J
  =\varepsilon J(A+\varepsilon I)^{-1},
  \label{eq:offtangent}
\end{equation}
and the last two terms of Eq.~\eqref{eq:linearized-expansion} are generally nonzero.
The middle term couples the learned increment to the fixed reconstruction residual;
it can shift the increment and is not a shrinkage penalty. The final term is an
anisotropic shrinkage. Diagonalizing $A$ with eigenvalues $\lambda_i$ and writing
$\delta=\sum_i\delta_i v_i$ gives
\begin{equation}
  \lVert B_\varepsilon\delta\rVert^2
  = \sum_{i=1}^{d} \frac{\lambda_i\,\varepsilon^2}{(\lambda_i+\varepsilon)^2}\,
  \delta_i^2
  \;\leq\; \frac{\varepsilon}{4}\,\lVert\delta\rVert^2 ,
  \label{eq:spectral}
\end{equation}
uniformly in $J$ and with no rank or conditioning assumption, since the scalar
coefficient is maximized at $\lambda_i = \varepsilon$, where it equals
$\varepsilon/4$, and vanishes both as $\lambda_i \to 0$ and as
$\lambda_i \to \infty$. With $w=0.3$ and $\varepsilon=10^{-5}$, the quadratic
component alone is bounded by
$7.5\times10^{-7}\lVert\delta\rVert^2$.

This bound does not control the middle cross-term. Equivalently, if $s_i$ denotes a
singular value of $J$, then $B_\varepsilon$ has singular values
$\varepsilon s_i/(s_i^2+\varepsilon)$, maximized at $s_i=\sqrt{\varepsilon}$, so
$\lVert B_\varepsilon\rVert_2\leq\sqrt{\varepsilon}/2$ and by Cauchy--Schwarz the
cross-term is at most
$w\sqrt{\varepsilon}\,\lVert\delta\rVert\,\lVert Q_\varepsilon r_t\rVert$ in absolute
value. The two surviving terms therefore scale as $\sqrt{\varepsilon}$ and
$\varepsilon$ respectively, and at $\varepsilon=10^{-5}$ they differ by roughly four
orders of magnitude: the cross-term, not the shrinkage term, is what the penalty
principally contributes.

We can bound its size from quantities we report. Treating the reconstruction residual
as approximately isotropic, a per-tenor RMSE of $0.659$ over $N{=}11$ tenors with a
five-dimensional tangent space gives
$\lVert Q_\varepsilon r_t\rVert\approx0.659\sqrt{11}\sqrt{6/11}\approx1.6$, so with
$w=0.3$ and $\sqrt{\varepsilon}\approx3.2\times10^{-3}$ the cross-term is bounded by
roughly $1.5\times10^{-3}\lVert\delta\rVert$ in the vector-norm convention of
Eq.~\eqref{eq:penalty-actual}, or $1.4\times10^{-4}\lVert\delta\rVert$ under the
mean-over-tenors reduction the implementation uses, against a fit loss of order
$10^{-2}$. The one input we do not have is the typical size of the one-step latent
increment $\lVert\delta_\psi\rVert$, which was not logged; for
$\lVert\delta_\psi\rVert$ of order $10^{-2}$ the bound sits three to four orders of
magnitude below $\mathcal{L}_{\mathrm{fit}}$. This is a bound rather than a
measurement, and the isotropy assumption is a convenience, but it fixes the scale of
the surviving linear signal as small without asserting that it is zero.
\paragraph{Decoder-curvature signal.} The implemented move is nonlinear rather than
$J\delta$. Locally,
\begin{equation}
  m_\psi=J\delta_\psi+\tfrac12 H[\delta_\psi,\delta_\psi]
  +O(\lVert\delta_\psi\rVert^3),
\end{equation}
where $H$ is the decoder Hessian at $z_t$. Under the exact projector,
\begin{equation}
  Q_0(r_t+m_\psi)
  =Q_0r_t+\tfrac12Q_0H[\delta_\psi,\delta_\psi]
  +O(\lVert\delta_\psi\rVert^3).
\end{equation}
When $Q_0r_t\ne0$, the cross-term between reconstruction error and curvature changes
the loss at order $\lVert\delta_\psi\rVert^2$ and its gradient at order
$\lVert\delta_\psi\rVert$. If $Q_0r_t=0$, the curvature contribution instead enters
the loss at fourth order. The actual nonlinear penalty can therefore train through
decoder curvature even without ridge leakage.
The analytic conclusion is consequently narrower than total inertness: the exact
local tangent projector removes the leading-order decoded move, while the
implemented objective retains a reconstruction-residual cross-term, ridge shrinkage,
and decoder-curvature effects. Our one-run ablation finds only small changes at the
trained horizons, which is consistent with these surviving signals being weak in
that run but does not prove that they are negligible. Establishing their scale would
require the penalty-gradient norm, $Q_\varepsilon r_t$, and the spectrum of
$J^{\top}J$ during training.
A nonlinear re-encoding construction based on $D_\theta(E_\phi(\hat y))$ would
target return to the decoder image rather than to a local affine tangent space; the
benchmark of \citet{NoArbVAEYieldCurves} uses such a re-encoding step. Another
nondegenerate design would apply tangent consistency to a curve-space proposal whose
first-order motion is not already generated through the decoder Jacobian.
\subsection{Large re-encoding error complicates projection}
\label{sec:analysis-reencoding}
Stage~A test reconstruction RMSE is $0.659$, so real curves sit far from
$D(E(y))$. A future method that maps forecasts onto $D(E(\hat y))$ need not move them
closer to the true curve $y^\star$: $D(E(\cdot))$ is not a metric projection, and
when re-encoding error is this large, mapping onto the decoder image can increase or
decrease forecast error. The observed error therefore cautions against assuming an
accuracy gain from projection, but it is not a lower bound on projected forecast
RMSE or an upper bound on possible improvement.
It also helps explain why absolute re-encoding distance stays near $0.65$ across
configurations: the metric closely tracks the Stage~A error at the persistence
anchor, so it is a weak discriminator. The re-encoding gain subtracts that shared
anchor value, but on our single run it does not order the two drift variants
consistently across horizons.
\subsection{Tangent residual: what it does and does not show}
\label{sec:analysis-tangent}
The tangent-move residual measures whether nonlinear decoded moves stay in the
column space of $J(z_t)$. Unlike absolute re-encoding distance it is not
saturated by Stage~A error, so it isolates the dynamics contribution. It is not the
quantity the training penalty evaluates: the penalty acts on the persistence-residual
forecast, where the parameter-dependent part is degenerate by
Section~\ref{sec:analysis-inert} at leading order under exact projection, while the
diagnostic acts on the decoded move alone.
The measured differences do not separate the two drift configurations.
\texttt{sde\allowbreak\_jacobian} is nominally lower at $h{=}1$, $5$, and $21$,
but by margins
($2.8$ versus $2.9\times10^{-5}$ at $h{=}1$) that we cannot attribute to the penalty
from one run. The metric does separate the regularizer bundle sharply---$0.0033$ at
$h{=}1$ rising to $0.85$ at $h{=}21$---and that contrast is where its diagnostic
value lies.
\subsection{The regularizer bundle as an unstable comparator}
\label{sec:analysis-pde}
The configuration carrying the bundle provides the sharpest contrast in the study. At
$h{=}1$ it already raises RMSE from $\approx0.029$ to $\approx0.041$; by $h{=}5$ RMSE
is $\approx0.197$; by $h{=}21$ it exceeds $1.75$; and at $h{=}63$ it exceeds $7.6$.
Tangent residuals grow in parallel.
The conclusion concerns the bundle as a whole, and four candidate causes remain
unseparated by this design. \emph{Horizon gating:} the bundle is applied at $h{=}1$
only, while the curve-fitting loss uses all three trained horizons, so a model tuned
to satisfy it at one step is free to drift at longer ones. \emph{Bundling:} the
discount-monotonicity and forward-smoothness terms accompany the residual and either
could drive the divergence. \emph{Unit and time-scale misspecification:} the
residual is computed on robust-scaled yields treated as if they were continuously
compounded rates, while a daily drift is combined with annual maturity units
(Section~\ref{subsec:constraints}). \emph{Structural misspecification:} the residual
omits both the Hessian convexity term and the $\tau\,\partial_\tau y$ term of
$\partial_\tau P$.
We therefore make no claim about no-arbitrage penalization as such, and none about
what a correctly specified residual would do. What the evidence supports is narrower:
this $h{=}1$-gated regularizer bundle is associated with RMSE growth by a factor of
roughly $40$ from $h{=}1$ to $h{=}21$ and by more than two orders of magnitude by
$h{=}63$, whereas configurations without it remain stable across all evaluated
horizons. A residual-only ablation in correct units, at matched horizons, is the
natural next experiment.
\subsection{Extrapolation to 63 business days}
\label{sec:analysis-h63}
At $h{=}63$, outside the training horizon set, persistence is best and
\texttt{sde\_only} has lower error than \texttt{sde\_jacobian} by $9.8\times10^{-3}$
--- two orders of magnitude larger than the differences at trained horizons, and in
the opposite direction. We treat this as extrapolation stress rather than evidence
about the penalty, and restrict our claims to trained horizons.
\subsection{Limitations}
\label{sec:analysis-limitations}
\begin{itemize}
  \item \textbf{Single run, no inference.} Every number is a point estimate from one
  seed with nondeterministic kernels. We performed no paired significance testing.
  Differences of order $10^{-5}$ between drift configurations, and the $1.3\%$ margin
  over persistence at $h{=}21$, are descriptive. A Diebold--Mariano test with HAC
  standard errors is the appropriate instrument, and the overlapping forecast windows
  at $h>1$ induce serial correlation such a test would need to accommodate.
  \item \textbf{Data pipeline defect.} Roughly $16\%$ of the panel consists of linear
  interpolants spanning the 2002--2006 discontinuation of the 30-year series
  (Section~\ref{sec:data}), plus interpolated market holidays. This window lies
  entirely within the training split, so it does not create synthetic rows in the
  held-out period; however, it affects the representation and scaler used to produce
  every held-out forecast. Stage~A was fit on partly synthetic data, and a rebuilt
  panel is the correct fix.
  \item \textbf{The regularizer is not a no-arbitrage penalty.} It is computed on
  robust-scaled par yields treated as continuously compounded zero rates, combines
  daily and annual time units, and omits two terms of the pricing PDE. Our
  instability finding is about that penalty as implemented, not about no-arbitrage
  penalization.
  \item \textbf{The ablation is not isolated.} The bundle carries three terms applied
  at a single horizon; the instability cannot be attributed to any one of them.
  \item \textbf{Classical baselines excluded.} PCA--VAR and NSS results are omitted
  because the rolling refit duplicated observations. Persistence is therefore the
  only valid external benchmark in this report.
  \item \textbf{Surviving penalty terms not measured.} The cross-term and
  curvature signals of Section~\ref{sec:analysis-inert} are bounded analytically but
  were not logged during training; the size estimate given there rests on an isotropy
  assumption and on an unmeasured typical increment norm.
  \item \textbf{Diffusion is not trained without the bundle.} No term of the
  \texttt{sde\_only} or \texttt{sde\_jacobian} objective depends on $\sigma$, so
  those configurations are drift networks rather than fitted SDEs.
  \item \textbf{Stage~A fidelity.} Reconstruction RMSE of $0.659$ makes an accuracy
  gain from re-encoding uncertain; it does not impose a formal bound on projected
  forecast error.
  \item \textbf{Soft penalties only.} Hard projection at inference was not evaluated.
  \item \textbf{Single market panel.} Results are for U.S.\ Treasury CMT yields;
  generalization to other curves, currencies, or crisis regimes is untested.
\end{itemize}
\section{Conclusion}
\label{sec:conclusion}
We set out to characterize a decoder-Jacobian tangent penalty for learned latent
yield-curve dynamics, on a shared frozen representation with a random walk as the
null. The main result is a leading-order degeneracy, not total inertness.
An exact tangent projector annihilates a first-order decoded latent move, making the
linearized loss constant in the dynamics parameters. The implemented objective,
however, uses a finite ridge and a nonlinear decoder. Its loss contains a cross-term
between the increment and Stage~A reconstruction residual, an anisotropic shrinkage
term, and a decoder-curvature contribution. The uniform $\varepsilon/4$ bound applies
only to the quadratic shrinkage term; the cross-term scales as $\sqrt{\varepsilon}$
and dominates it. The actual penalty can therefore train, though
not through the direct first-order tangent signal its geometric interpretation might
suggest. This decomposition is provable without reference to any run. Empirically,
switching on the penalty changes trained-horizon RMSE by at most
$7\times10^{-5}$ in our single run; this is consistent with weak surviving signals
in that run but does not establish that they are always negligible.
The remaining findings are descriptive and hold for a single run. Persistence is the
binding benchmark at short horizons and beyond the trained horizon set; the drift
model improves on it at $h{=}21$ by roughly $1.3\%$ in point estimate; and an
$h{=}1$-gated bundle of curve-shape regularizers destabilizes multi-step forecasts by
two orders of magnitude, though the design cannot say which of its terms or
misspecifications is responsible.
The methodological lesson runs through all of it: a constraint ablation is only
interpretable against a random-walk null, with each penalty isolated, in units where
the penalized quantity means what its name says, and with enough runs to know which
differences are real. The natural next steps are a residual-only ablation in correct
units at matched horizons, a re-encoding form of the geometric penalty, a multi-seed
protocol with paired tests, and a Stage~A model whose reconstruction error is low
enough that re-encoded points approximate market curves. Those steps would determine
whether a nondegenerate geometric penalty improves either forecast accuracy or
curve-consistency diagnostics.
\label{last_page}
\bibliography{references}
\bibliographystyle{iclr2022_conference}
\end{document}
